\section{Kaleidoscopes and Mirrors: Coxeter Groups}

Where do these perfect tilings come from? It is not someone carefully laying down tile after tile. There is a far more elegant answer. The whole infinite tiling is generated automatically, from a single tile, by mirrors. The mathematics of this is the theory of Coxeter groups, named after the geometer Harold Coxeter. You do not need the mathematics. You need the kaleidoscope.

\subsection{The kaleidoscope}

You have looked into a kaleidoscope. A few bits of colored glass sit between mirrors, and you see a gorgeous, endless, perfectly symmetric pattern. The bits of glass are few. The pattern is vast. The mirrors did all the work, reflecting the scrap of color over and over until it filled the view.

That is exactly how a regular tiling is made. Take one tile. Stand mirrors along its edges. Each mirror reflects the tile to make a copy on the other side. Those copies have edges too, with their own mirrors, making more copies. Reflections of reflections of reflections, on and on, and the whole space fills up with identical tiles. You never placed them. The mirrors generated them.

A Coxeter group is just the bookkeeping for all those reflections: every way you can combine ``flip across this mirror, then that one, then that one'' to land on some copy of the original tile. The group is the set of all those moves. The tiling is what you see when you apply all of them to one tile.

\subsection{The only thing you have to choose}

Here is the beautiful part, the part that matters for the whole theory. To specify a kaleidoscope, you do not need much. You only need the angles between the mirrors. And for these regular tilings, those angles are set by whole numbers.

That is it. A couple of integers fix the angles, the angles fix the mirrors, the mirrors generate the tiling, and the tiling fixes the curvature of the space. Everything cascades from a tiny handful of whole numbers. Nothing is tuned by hand. Nothing is arbitrary. You pick two or three integers, and an entire infinite, perfectly regular, curved crystal pops out, uniquely determined.

This is why the theory can claim its foundation has nothing arbitrary in it. The crystal is not a design choice with a thousand knobs. It is the automatic output of a kaleidoscope whose only settings are a few integers.

\subsection{From symmetry to crystal}

There is one more layer worth noticing, because it connects to the bottom of the whole theory. Reflections are a kind of symmetry, a way of moving a shape so it lands on itself. The integers that set the mirror angles are, deep down, the symmetries of the whole numbers themselves. So the chain runs: the plain integers, then their symmetries, then those symmetries widened into mirrors, then the mirrors generating a hyperbolic crystal. Arithmetic becomes geometry. We will come back to this ladder near the end of Part II. For now, just hold the headline.

\subsection{What you now know}

You know that the crystal is not built tile by tile. It is generated by reflection, like a kaleidoscope, from one tile and a few mirror angles. And those angles are nothing but whole numbers. The next chapter gives those numbers their proper names, so we can say precisely which crystal we mean.

\textbf{Takeaway:} a regular tiling is generated automatically from one tile by mirrors, like a kaleidoscope, and the only settings are a few whole numbers (the mirror angles). A tiny handful of integers determines an entire infinite curved crystal, with nothing arbitrary added.
